
\chapter*{Part I Integrated Lab: Security Model Review}
\addcontentsline{toc}{chapter}{Part I Integrated Lab: Security Model Review}

\section*{Purpose}

This Integrated Lab replaces a separate capstone. It is the practical checkpoint for Part I. The goal is to combine the first seven chapters into one professional security-model review package.

Part I taught the reusable method:

\begin{lstlisting}
model the state machine
map assets and authority
inspect transactions and authorization
identify commitments, proofs, and trust boundaries
write a threat model
write invariants
reason about exploitability
\end{lstlisting}

This lab asks you to perform that method on one small protocol.

\section*{Scenario}

You are reviewing a toy yield vault called \passthrough{\lstinline!EpochVault!}.

\begin{lstlisting}
Users deposit USDC into the vault.
The vault mints shares.
A strategist can allocate USDC into approved strategies.
An oracle reports the strategy value once per epoch.
Users request withdrawals into a queue.
Withdrawals are processed at the next epoch boundary.
Governance can add strategies after a 48-hour timelock.
A guardian can pause deposits immediately.
A separate emergency role can pause withdrawals.
The vault is upgradeable through a proxy admin controlled by governance.
\end{lstlisting}

Assume the protocol documentation claims:

\begin{quote}
Users own vault shares backed by USDC and approved strategies. Governance protects users through timelocked upgrades. The oracle reports fair strategy value. The guardian can pause during emergencies.
\end{quote}

Your job is to replace that product description with a security model.

\section*{Required Output}

Produce one document titled:

\begin{lstlisting}
Security Model Review Package: EpochVault
\end{lstlisting}

It must include the sections below.

\section*{1. State-Machine Model}

Include:

\begin{itemize}
\tightlist
\item
  state variables;
\item
  accepted inputs;
\item
  transition table;
\item
  rejection conditions;
\item
  forbidden states;
\item
  at least three safety properties;
\item
  at least two liveness properties.
\end{itemize}

\section*{2. Asset Inventory}

Identify assets that are not merely token balances.

At minimum include:

\begin{itemize}
\tightlist
\item
  USDC deposits;
\item
  vault shares;
\item
  strategy positions;
\item
  queued withdrawal claims;
\item
  oracle-reported value;
\item
  governance power;
\item
  guardian and emergency roles;
\item
  upgrade authority.
\end{itemize}

\section*{3. Authority Map}

Create a table with:

\begin{lstlisting}
asset or state
transition
required authority
authority holder
delay or timelock
observability
worst-case failure
recovery path
\end{lstlisting}

Then answer:

\begin{lstlisting}
Which authority is economically equivalent to control over the vault?
Which authority can break liveness without directly stealing funds?
Which authority can change the future transition rules?
\end{lstlisting}

\section*{4. Trust-Boundary Map}

Draw or describe the trust boundaries among:

\begin{itemize}
\tightlist
\item
  users and wallets;
\item
  frontend or transaction builder;
\item
  vault contract;
\item
  USDC token contract;
\item
  strategy contracts;
\item
  oracle or reporter;
\item
  governance and timelock;
\item
  guardian and emergency role;
\item
  proxy admin or upgrade path.
\end{itemize}

For each boundary, state what is trusted, who can lie or fail, who verifies, and what happens if the assumption is false.

\section*{5. Transaction and Authorization Review}

Choose three transitions:

\begin{itemize}
\tightlist
\item
  deposit;
\item
  request withdrawal;
\item
  governance strategy addition or vault upgrade.
\end{itemize}

For each transition, identify:

\begin{itemize}
\tightlist
\item
  caller or signer;
\item
  authorization rule;
\item
  nonce or replay protection if relevant;
\item
  state read;
\item
  state written;
\item
  external calls;
\item
  events or logs;
\item
  valid-but-unsafe outcome.
\end{itemize}

\section*{6. Threat Model}

Produce a one-page threat model with:

\begin{itemize}
\tightlist
\item
  protocol goal;
\item
  actors;
\item
  assets;
\item
  critical transitions;
\item
  privileged roles;
\item
  external dependencies;
\item
  top safety assumptions;
\item
  top liveness assumptions;
\item
  top five attack hypotheses;
\item
  review priorities.
\end{itemize}

\section*{7. Invariant Catalogue}

Write at least eight invariants:

\begin{itemize}
\tightlist
\item
  two conservation invariants;
\item
  two authorization invariants;
\item
  two share/accounting invariants;
\item
  one temporal or withdrawal invariant;
\item
  one recovery or pause invariant.
\end{itemize}

Each invariant must state its assumptions and what evidence would falsify it.

\section*{8. Exploitability Memo}

Pick one attack hypothesis and write an exploitability memo with:

\begin{itemize}
\tightlist
\item
  broken invariant;
\item
  preconditions;
\item
  required capital;
\item
  required timing;
\item
  liquidity assumptions;
\item
  ordering assumptions;
\item
  expected profit or damage;
\item
  reliability;
\item
  mitigations;
\item
  residual risk.
\end{itemize}

\section*{Grading Criteria}

The submission is acceptable if it:

\begin{itemize}
\tightlist
\item
  identifies state and transitions before naming bug classes;
\item
  distinguishes ownership, custody, authority, and economic exposure;
\item
  names trust assumptions instead of hiding them behind words like ``trusted'' or ``decentralized'';
\item
  writes invariants that are specific, falsifiable, and linked to impact;
\item
  explains exploitability using capital, timing, liquidity, ordering, reliability, and profit;
\item
  produces artifacts another reviewer could use.
\end{itemize}
